\begin{theoremEnd}[combined main,category=error]{proposition}[Posterior error from a conservative continuum threshold]
\CombinedResultCheck{prop:error}{proposition}\label{prop:error}
Under \cref{thm:finite}, let $\widetilde\tau_r$ be measurable and satisfy
$\tau_r\leq\widetilde\tau_r\leq\tau_r+\delta$ almost everywhere
under the ordinary posterior of $\boldsymbol g$, for a fixed $\delta\geq0$.
Use the same upper threshold in marginal weighting and reconstruction.
Let $P_{\boldsymbol y}$ be the ordinary coefficient posterior, and use
$\mathcal C_r$ also for its coefficient preimage. Define
$\widetilde{\mathcal C}_r=\{\gamma\geq\widetilde\tau_r(\boldsymbol g)\}$.
Both conditioned laws exist and, with $\TV(P,Q)=\sup_A|P(A)-Q(A)|$,
\begin{equation}\label{eq:tv}
 \begin{split}
 \TV\bigl(P_{\boldsymbol y}(\cdot\mid\widetilde{\mathcal C}_r),
          P_{\boldsymbol y}(\cdot\mid\mathcal C_r)\bigr)
 &=1-\frac{P_{\boldsymbol y}(\widetilde{\mathcal C}_r)}
                 {P_{\boldsymbol y}(\mathcal C_r)}
 \leq\min\left(1,\frac{\delta}
 {s\sqrt{2\pi}P_{\boldsymbol y}(\mathcal C_r)}\right).
 \end{split}
\end{equation}
Every measurable summary of the reconstructed coefficients inherits
this bound. The identity holds for any inner event of positive probability.
\end{theoremEnd}
\begin{proofEnd}
\CombinedProofCheck{prop:error}
Fix the ordinary Gaussian law, selected order, and almost-everywhere
error condition of \cref{prop:error}, modifying the numerical threshold
on the exceptional null set if necessary. Finite thresholds and $s>0$
give positive conditional tails, so both conditioning events have
positive probability. To compare the resulting laws, put
$A=\mathcal C_r$ and $B=\widetilde{\mathcal C}_r$. If the omitted
mass is zero, both conditioned laws agree. Otherwise the law on $A$
is a mixture of the laws on $B$ and $A\setminus B$, with weights
$P_{\boldsymbol y}(B)/P_{\boldsymbol y}(A)$ and its complement.
The components have disjoint supports, so total variation equals the
second weight, establishing the identity for any positive-probability
inner event.

For the threshold events considered here, conditional
Gaussian integration identifies the omitted mass as
\[
 \E\left[\Phi((\widetilde\tau_r(\boldsymbol g)-\mu(\boldsymbol g))/s)
          -\Phi((\tau_r(\boldsymbol g)-\mu(\boldsymbol g))/s)
          \mid\boldsymbol y\right].
\]
Since the conditional density is at most $1/(s\sqrt{2\pi})$ and the
excluded interval has length at most $\delta$, dividing this mass
bound by the probability of the exact event proves \cref{eq:tv}.
The invertible coefficient map preserves total variation, while
measurable summaries cannot increase it, proving the remaining assertion.
\end{proofEnd}
